\documentclass[10pt,a4paper]{amsart}
\usepackage{amsmath,amssymb,amsthm,graphicx}
\usepackage{amsmath}
\usepackage{amssymb}
\usepackage{hyperref}
\newtheorem{Theorem}{Theorem}[section]
\newtheorem{Corollary}[Theorem]{Corollary}
\newtheorem{Lemma}[Theorem]{Lemma}
\newtheorem{Proposition}[Theorem]{Proposition}
{ \theoremstyle{definition}
	\newtheorem{Definition}[Theorem]{Definition}
	\newtheorem{Note}[Theorem]{Note}
	\newtheorem{Example}[Theorem]{Example}
	\newtheorem{Remark}[Theorem]{Remark} }

\usepackage[all]{xy}


\def\KK{\ensuremath\mathbb{K}}
\def\MM{\ensuremath\mathbb{M}}
\def\BB{\ensuremath\mathbb{B}}
\def\GL{\ensuremath\mathbf{GL}}
\def\UU{\ensuremath\mathbf{U}}
\def\PP{\ensuremath\mathbb{P}}
\def\CC{\ensuremath{\mathbb{C}}}
\def\NN{\ensuremath{\mathbb{N}}}
\def\QQ{\ensuremath{\mathbb{Q}}}
\def\RR{\ensuremath{\mathbb{R}}}
\def\ZZ{\ensuremath{\mathbb{Z}}}

\newcommand{\profinQ}{\mathfrak{Q}}
\newcommand{\reals}{\mathbb{R}}
\newcommand{\complexs}{\mathbb{C}}
\newcommand{\naturals}{\mathbb{N}}
\newcommand{\integers}{\mathbb{Z}}
\newcommand{\rationals}{\mathbb{Q}}
\newcommand{\quaternions}{\mathbb{H}}
\newcommand{\hyperbolic}{\mathbb{H}}
\newcommand{\K}{\mathbb{K}}
\newcommand{\realsorcomplexs}{\mathbb{K}}
\newcommand{\finitefield}{\mathbb{F}}
\newcommand{\Proj}{\mathbb{P}}
\DeclareMathOperator{\id}{id}
\DeclareMathOperator{\Ind}{Ind}
\newcommand{\boundary}[1]{\partial#1}
\newcommand{\boundedops}{\mathcal{B}}
\newcommand{\bundleover}{\!\!\downarrow\!\!}
\newcommand{\abs}[1]{\left\lvert#1\right\rvert} %absolute value
%\newcommand{\bigabs}[1]{\left\lvert#1\right\rvert} %absolute value
\newcommand{\norm}[1]{\left\lVert#1\right\rVert}
\newcommand{\generate}[1]{\langle#1\rangle}
\newcommand{\tensor}{\otimes}
\newcommand{\into}{\hookrightarrow}
\newcommand{\onto}{\twoheadrightarrow}
\newcommand{\iso}{\cong}
\newcommand{\disjointunion}{\amalg}
\newcommand{\subgroup}{\leq}
\newcommand{\supergroup}{\geq}
\newcommand{\superset}{\supset}
\newcommand{\normalsubgroup}{\lhd}
\newcommand{\semiProd}{\rtimes}
\newcommand{\semiprod}{\semiProd}
\newcommand{\dividesnot}{\nmid}
\newcommand{\divides}{\mid}
\DeclareMathOperator{\Pos}{Pos}
\DeclareMathOperator{\Aut}{Aut}
\DeclareMathOperator{\supp}{supp} %support
\DeclareMathOperator{\clos}{clos} %closure
\DeclareMathOperator{\im}{im} %image
\DeclareMathOperator{\vol}{vol} %Volume
\DeclareMathOperator{\diam}{diam} %Diameter
\DeclareMathOperator{\dist}{dist} %Distance
\DeclareMathOperator{\ord}{ord} %order
\DeclareMathOperator{\End}{End} %Endomorphisms
\DeclareMathOperator{\Hom}{Hom} %Homomorphisms
\DeclareMathOperator{\Span}{span}
\DeclareMathOperator{\spin}{spin}
\DeclareMathOperator{\diag}{diag}
\DeclareMathOperator{\spec}{Spec}
\DeclareMathOperator{\trace}{trace}
\DeclareMathOperator{\tr}{tr}
\DeclareMathOperator{\pr}{pr}
%\DeclareMathOperator{\res}{res} %Restriction
\DeclareMathOperator{\grad}{grad} %Restriction
% \DeclareMathOperator{\Sp}{Sp} %Spur % schon definiert???
\DeclareMathOperator{\fl}{f\/l} %flip
\DeclareMathOperator{\coker}{coker}
\DeclareMathOperator{\Char}{Char}
\DeclareMathOperator{\ind}{ind}
\DeclareMathOperator{\sgn}{sgn}
\DeclareMathOperator{\Ric}{Ric} % Ricci curvature
\DeclareMathOperator{\scal}{scal} % Scalar curvature
\DeclareMathOperator*{\dirlim}{dirlim}
\DeclareMathOperator*{\invlim}{invlim}
\DeclareMathOperator{\rank}{rank}
\DeclareMathOperator{\coind}{coind}
\DeclareMathOperator{\res}{res}
\newcommand{\bigdisjointunion}{\bigsqcup}
\DeclareMathOperator{\Orb}{Orb}
\DeclareMathOperator{\GTOP}{G-TOP}
\newcommand{\forget}[1]{}
\newenvironment{Forget}{}{}
\newcommand{\innerprod}[1]{\langle #1 \rangle}
\title{The classifying space of the equivariant fundamental groupoid is universal for equivariant cellular maps}
\author{Massimiliano Puglisi, Thomas Schick, Vito Felice Zenobi}
\date{doi:10.3842/SIGMA.2025.093; CC BY 4.0}
\begin{document}
\maketitle
\noindent\textit{Editorial note.} All delivered mathematical text is the author's own, copied from the pinned publisher source sources/geometry-fl04-stolz-positive-scalar-curvature-st/pinned-original-source.tex; no statement or proof step was rewritten, reordered or removed.\par\medskip
\noindent\textit{Editorial note.} Preamble provenance: the retained preamble is the source's own layout adapter, if any, plus the author's macro block at lines 11-108 of sources/geometry-fl04-stolz-positive-scalar-curvature-st/pinned-original-source.tex (everything between the publisher class metadata and \textbackslash{}begin\{document\}, comments excluded). Nothing was added, no macro was invented, and the publisher document class is deliberately not reproduced; the excerpt serves only as preamble, never as a span coordinate.\par\medskip
\noindent\textit{Editorial note.} Metadata provenance: source\_id and source\_version were taken from sources/geometry-fl04-stolz-positive-scalar-curvature-st/license-observation.json; the title from SOURCE.json; the author list from the source metadata record; the source URL from the first non-empty of source\_url, source\_urls[0] and the evidence's fixed\_version\_url/license\_url. The licence evidence file used by this card is sources/geometry-fl04-stolz-positive-scalar-curvature-st/license-observation.json, recorded in license.evidence\_path so that the binding target is explicit.\par\medskip
\noindent\textit{Editorial note.} External references are navigation pointers only; every deliberately kept step of the author's argument stays inside the counted spans. The displayed numbers of those pointers are the author's own numbering under the author's \textbackslash{}newtheorem declarations (section.index as declared), computed mechanically by the lane's scanner.\par\medskip
\noindent\textit{Editorial note.} The counted proof contains no reference to any labelled item and no bibliography citation, so the card needs neither a navigation stub nor a bibliography excerpt.\par\medskip
\noindent\textit{Editorial note.} The uniqueness half is described by the author as following the pattern of the non-equivariant case together with the same homotopy vanishing; that sentence is kept verbatim inside the span.\par\medskip
\noindent\textit{Allowed standard prerequisites (editorial).}\begin{itemize}
\item The equivariant fundamental groupoid Pi\_1(Y;G) and its classifying space B Pi, with isotropy subgroups H and the fixed point sets (B Pi)\textasciicircum{}H. Set up in the paragraphs preceding the counted statement, ending at source line 1068. Fixes the objects compared and the skeleta on which the map is constructed.
\item Vanishing of the higher homotopy groups of the fixed point sets of B Pi. Used inside the counted proof at lines 1114-1117 and 1134-1135. Makes every attaching map of a higher cell contractible in the relevant fixed point set, which is what allows the extension over the skeleta.
\end{itemize}\medskip
	\begin{Proposition}\label{universalspace}
		For every $G$-CW-complex $Y$ and for every natural transformation $\Phi\colon
		\Pi_1(Y;\allowbreak G)\to \Pi$, there exists, unique up to $G$-equivariant
		homotopy, a $G$-equivariant cellular map
		$\varphi\colon Y\to B \Pi$
		such that
		$\varphi_{\#}=\Phi$.
	\end{Proposition}

	\begin{proof}
		We begin by defining the map $\varphi$ on the $G$-equivariant $0$-skeleton of $Y$ by putting
		$\varphi_{|Y_{(0)}}:=\smash{\Phi(\{e\})_{|Y_{(0)}}}$.
		Now, we proceed to define $\varphi$ on the $1$-skeleton of $Y$. For a $G$
		$1$-cell
		$c$ of the form $G/H\times [0,1]$ (with isotropy $H$) pick the ordinary
		$1$-cell $c_H:=H/H\times D^1$ contained in $Y^H$. Then $c_H$
		defines (if we choose an orientation of it) an element \smash{$\gamma \in
		\Pi_1\bigl(Y^H\bigr)_{|Y^H_{(0)}}$}. Define now \smash{$\phi\colon c_H\to \bigl(B\Pi^H\bigr)^{(1)}$} such that
		this map represents the element $\Phi(H)(\gamma)$ in
		$\Pi_1\bigl(B\Pi^H\bigr)|_{ (B\Pi^{ (0)} )^H}$. Since $H$ acts trivially on the target, this extends
		uniquely to a $G$-map on the $G$-cell $c$ with values in the
		$1$-skeleton of $B\Pi$. Of course, we could have picked a different base cell
		in the $G$-cell $c$. But because $\Phi$ is a natural transformation, the
		resulting map is independent of this ---up to the choice of the representative in the
		homotopy class of $\Phi(H)(\gamma)$ which had already to be made anyway.
		This defines the map $\phi$ on the $1$-skeleton of $Y$.
		
		To extend $\phi$ to $Y^{(2)}$, let $c$ be a $G$-$2$-cell of the form
		$G/H\times D^2$. Pick the corresponding ordinary $2$-cell $H/H\times
		D^2$ (with isotropy $H$)
		which is contained in $Y^{H}$. Its attaching map~${\psi\colon S^1\to Y^{H}}$
		is obviously contractible in $Y^{H}$ and hence trivial (i.e., a unit) in
		the fundamental groupoid~\smash{$\Pi_1\bigl(B\Pi^H\bigr)|_{ (B\Pi^{ (0)})^H}$} (we conjugate with a path in
		$Y^H$ to the $0$-skeleton to make~$\psi$
		represent an element of \smash{$\Pi_1\bigl(Y^H\bigr)_{B\Pi^H_{(0)}}$}, the triviality does not
		depend on the choice of such a path). Consequently, $\Phi(H)$ being a map of
		groupoids, also the image under $\Phi(H)$ of this element of the fundamental
		groupoid is trivial (a unit). By construction of $\phi$ on the $1$-skeleton,
		this image element is represented by $\phi\circ\psi$ (up to the chosen
		conjugation with a path to the $0$-skeleton), which is hence contractible in
		$B\Pi^H$. Extend $\phi$ over $H/H\times D^2$ using this contraction and then
		extend it $G$-equivariantly to $c=G/H\times D^2$.
		
		Inductively, we then extend $\phi$ over the $k$-skeleta of $Y$. The extension
		property now follows because each attaching map has contractible image by the
		vanishing of all higher homotopy groups of all components of all fixed point
		sets of $B\Pi$ for the various subgroups of $H$.
		
		By the very construction of $\phi$ on the $1$-skeleton, we have
		$\varphi_{\#}=\Phi$, because the morphism sets of
		\smash{$\Pi(H)=\Pi_1\bigl(B\Pi^H\bigr)|_{ (B\Pi^{ (0)})^H}$} are
		generated by the $1$-cells.
		
		The last step of the proof concerns uniqueness. Choose therefore a
		$G$-equivariant map
		$\varphi'\colon Y \to B \Pi$ such that
		$\varphi'_{\#}=\Phi$. We have to show that $\varphi'$ is $G$-homotopic to
		$\varphi$.
		
		It is immediate from the definition that if $\varphi'_{\#}=\varphi_{\#}$, then
		their restrictions to $Y_{(0)}$ are equal. The construction of the desired
		$G$-homotopy is now done inductively over the skeleta and follows the pattern
		in the non-equivariant case and for the construction part, making use of the
		condition that $\phi_\#'=\phi_\#$ for the extension over the $1$-skeleton and
		of the vanishing of higher homotopy groups for the further extensions.
	\end{proof}

\end{document}
